\scanpage{10}
\mspage{\unsure{9|6}}
\mspage{34}
\[
(D-\operatorname{div}(g)).(w_1\times\cdots\times A\times\cdots\times w_5)=0.
\]
($h$-ième facteur égal à $A$ dans la seconde parenthèse). Supposons le diviseur $D-\operatorname{div}(g)$ non nul, et soit $Z$ l'une de ses composantes. D'après la relation précédente, $Z$ se projette suivant une sous-variété de codimension $\geq2$ sur chacun des produits partiels formés avec quatre des facteurs du produit $B=A\times A\times A\times A\times A$. Ceci est impossible, puisque $Z$ est de codimension $1$ relativement à $B$. On a bien $D=\operatorname{div}(g)$.

Soit maintenant $z$ un point générique de $A$ sur le corps $K(a_0,a_1,a_2,b_0)$, et considérons la fonction $g'$ sur $A$, définie sur $K$, telle qu'on ait
\[
g'(z)=g(a_0,a_1,a_2,b_0,z).
\]
D'après ([14], IX, 2, th.3), on a
\[
a_0\times a_1\times a_2\times b_0\times\operatorname{div}(g')=\operatorname{div}(g).(a_0\times a_1\times a_2\times b_0\times A).
\]
Or, d'après l'expression de $\operatorname{div}(g)$ obtenue plus haut, et d'après le lemme 2, on a
\[
\operatorname{div}(g).(a_0\times a_1\times a_2\times b_0\times A)=a_0\times a_1\times a_2\times b_0\times X_{\underline{a}}.
\]
On a donc
\[
\operatorname{div}(g')=X_{\underline{a}}.
\]
De plus, le point $c$, ainsi que le point $c_0=a_0\times a_1\times a_2\times b_0\times b_0$, n'appartiennent pas au support de $\operatorname{div}(g)$. La fonction $g$ est donc définie en $c$ et en $c_0$, et on a les relations
\[
g(c)=g'(b_1)\qquad g(c_0)=g'(b_0).
\]
Remarquons en outre que, si $\bar{y}_1$ est un point générique de $A$ sur $K(x_0,x_1,x_2,y_0,y_1)$, on a
\[
g(x_0,x_1,x_2,y_1,\bar{y}_1)=g(x_0,x_1,x_2,y_0,\bar{y}_1)\,g(x_0,x_1,x_2,y_0,y_1)^{-1}.
\]
